\documentclass[11pt,oneside]{article}
\usepackage[margin=1in]{geometry}
\usepackage{amsmath,amssymb,amsthm,mathtools,booktabs}
\usepackage[colorlinks=true,linkcolor=blue,urlcolor=blue]{hyperref}
\newtheorem{theorem}{Theorem}[section]
\newtheorem{proposition}[theorem]{Proposition}
\newtheorem{record}[theorem]{Record}
\newtheorem{firewall}[theorem]{Claim firewall}
\DeclareMathOperator{\SU}{SU}
\DeclareMathOperator{\Tr}{Tr}
\title{Renewal Yang--Mills Mass-Gap Research Notes\\
Post-Fifteenth-Archive Compact Lineage, Version V3}
\author{Continuation with a scope-checked result ledger}
\date{Research continuation}
\begin{document}
\maketitle
\tableofcontents

\section{Context, restart, and the unchanged objective}
\label{f16v1:context}

The objective is a rigorous nontrivial interacting four-dimensional
continuum Yang--Mills theory with a positive physical Hamiltonian mass
gap. That objective remains unproved. We currently work with an exact
compact \(\SU(2)\) Wilson transfer, its gauge constraints, and particular
weak-field limits. Auxiliary Gaussian separation, a lattice transfer
ratio, and a calibrated physical continuum mass are different claims.

The hundred-version fifteenth lineage is preserved in
\begin{center}\small
\path{Renewal_Yang_Mills_Mass_Gap_Research_Notes_V100_f15.tex}.
\end{center}
This is a compact restart, not a replacement of that archive.
References such as f15 V99 specify sections of that file.
Prior hypotheses and corrections remain binding. The summary below
is a result map, not a fresh certification of every historical claim.

The required V100 checkpoint reread SM continuum V72, NS stability V26,
shell-mixing V16, parallel YM V5, SM source-selection V31,
regenerative stationarity V34 and exact-action Einstein bridge V24.
Their current versions and hashes were unchanged. No companion
provides the missing volume-uniform native YM spectral comparison.
The useful cross-program lesson is to control densities or local
quantities rather than demand a volume-independent tail for an
extensive count. The detailed checkpoint is in f15 V100.
The next companion review is this lineage's V5; its next normal
archive/restart is V100.

The recent upstream change is to retain the normalized magnetic
image of the cube carrier rather than freeze its internal variables
or project its image straight back to the old ground band.
The present note continues by computing the exact overlap of those
conditional vectors across a native bridge kinetic step.
Its weak-field response is not an invariant-subspace theorem:
even the resulting Gaussian compressed family fails the semigroup
law at finite time. We keep that discrepancy explicit.

\section{Major interfaces retained from f15}
\label{f16v1:ledger}

\begin{center}\small
\begin{tabular}{@{}p{0.15\textwidth}p{0.79\textwidth}@{}}
\toprule
Range & Established objects and scope limits\\
\midrule
V1--V11 &
Full compact source/entropy estimates, literal blocking, growing-depth
curvature comparisons and scale-correct pinned-fibre interfaces.
They do not give physical continuum coercivity.\\[1mm]
V12--V37 &
Local chart and source-response comparisons with normalization,
exceptional-set and entropy costs retained; obstructions to dropping
those costs. No unqualified extensive-source mixing theorem.\\[1mm]
V38--V56 &
Native compact conditional landscapes, coherent boundary minima,
reflected source laws and paid pressure comparisons.
Rare or low-action regions are not automatically harmless for spectra.\\[1mm]
V57--V65 &
Positive Wilson transfer and vacuum response estimates in their
selected regimes; soft microscopic spectral limits.
A collapsing one-step lattice gap is not a disproof of a positive
gap after independently calibrated physical time scaling.\\[1mm]
V66--V81 &
Exact decimation/conditional laws, local Dirichlet and bad-region
control, full Gaussian invariant sectors, and vacuum-independent
transfer/compound-operator certificates.
A local gap or a finite matrix check is not the full native gap.\\[1mm]
V82--V91 &
Exact Gauss-boundary charge decomposition, magnetic charge walks,
charged-sector domination, compact endpoint channels and paid
spectral tails. Literature and undecidability audits import
methods, not an unverified mass-gap conclusion.\\[1mm]
V92--V97 &
Complete native coupled finite-region kernels and their growing
boundary carriers; full internal oscillator spectra and
finite-volume complement approximations.\\[1mm]
V98--V100 &
Actual spatial assembly, a sharp auxiliary coarse electric edge,
exact magnetic Gram/dressing, a fixed-volume Gaussian response,
and extensive old-oscillator excitation budgets.
Volume-uniform native dressed dynamics remain open.\\
\bottomrule
\end{tabular}
\end{center}

\subsection{The exact spatial assembly and its unresolved dynamics}

Let the spatial torus have side \(2m\), \(m\ge4\), with
\(B=m^3\) disjoint vertex cubes. A link is internal when its base
coordinate in its own direction is even; otherwise it is a bridge.
There are \(12B\) of each. Of the \(24B\) plaquettes, \(6B\)
are internal, \(12B\) cross one block interface and \(6B\) cross
in both directions. The coarse graph has four parallel bridges
between each pair of positive-direction neighbouring blocks.

A rooted spanning tree in each cube gives exact Haar coordinates
\(g\in G^{7B}\), \(X\in G^{5B}\), \(Y\in G^{12B}\), \(G=\SU(2)\).
After nonroot Gauss reduction the whole physical Hilbert space is
\begin{equation}
 \mathcal H_{\rm phys}
       =L^2(G^{5B}\times G^{12B})^{G^B}.
 \label{f16v1:physical}
\end{equation}
The root group conjugates its internal chords and acts at all its
bridge endpoints. Mixed plaquette words are exactly the two-bridge,
two-internal-edge words or the four-bridge words from the original
lattice. No plaquette is discarded.

For \(s(U)=1-\frac12\Tr U\), define the normalized central convolution
kernel and bridge operator
\begin{equation}
 k_\gamma(U)=\frac{e^{-\gamma s(U)}}{z_\gamma},\quad
 z_\gamma=\int_Ge^{-\gamma s(U)}\,dH(U),\quad
 \mathcal C_\gamma=\bigotimes_{e\ {\rm bridge}} C_\gamma^{(e)}.
 \label{f16v1:kinetic}
\end{equation}
These are the Wilson kinetic kernels, not exactly heat kernels.
The full native transfer has the exact factorization
\begin{equation}
 T_{\rm full,\gamma}
 =M_\times\left[
       \left(\bigotimes_b T_{\Box,\gamma}^{(b)}\right)
                       \otimes\mathcal C_\gamma\right]M_\times,
 \quad M_\times=e^{-\gamma S_\times/2},
 \label{f16v1:full}
\end{equation}
where \(S_\times\) sums all \(18B\) mixed plaquette actions.
The regional transfer factors and the bridge factor commute before
the magnetic sandwich, but the magnetic factor cannot be omitted.

The ground-carrier compression of the cut transfer in f15 V98
has coarse electric operator
\[
 \mathcal D=\mathcal E_{\rm bridge}
       +\sum_b\sum_{v,w}(L_\Box^+)_{vw}J_{bv}\cdot J_{bw},
 \qquad \mathcal E_{\rm bridge}\le\mathcal D
                         \le4\mathcal E_{\rm bridge}.
\]
On the coarse physical space its first nonzero eigenvalue is \(19/2\).
This is an auxiliary operator. Its proved native growing-carrier
error \(O_R(B/\sqrt\gamma)\) is larger than the \(1/\gamma\)
scale of that auxiliary heat edge. The magnetic image contributes
an additional order-one effect. No full-transfer gap follows.

\subsection{The exact normalized carrier and the leading response}

Here are sufficient definitions to continue without reprinting the
finite-region proofs. In one cube let \(\mathsf B\) be the
root-deleted edge incidence matrix, with row \(z_t-z_s\);
let \(\mathsf J\) inject its five chord coordinates; and let
\(\mathsf Z\mathsf B=0,\ \mathsf Z\mathsf J=I\) be the cycle matrix.
Use the rooted axial tree of f15 V97. Put
\[
 \mathsf F=-(\mathsf B^{\mathsf T}\mathsf B)^{-1}
                     \mathsf B^{\mathsf T}\mathsf J,\quad
 \mathsf M=\mathsf Z\mathsf Z^{\mathsf T},\quad
 U=\mathsf J+\mathsf B\mathsf F=\mathsf Z^{\mathsf T}\mathsf M^{-1}.
\]
Let \(\mathsf H\) be the Gram of the six linearized cube faces and
\(T=\mathsf M^{1/2}\mathsf H\mathsf M^{1/2}\).
The spectrum of \(T\) is \(4,4,4,6,6\).
The covariance of one colour component of the squared internal
Gaussian ground vector is
\begin{equation}
 C=\mathsf M^{1/2}[2\sqrt{T+T^2/4}]^{-1}\mathsf M^{1/2},
 \qquad C_B=\bigoplus_b C.
 \label{f16v1:covariance}
\end{equation}
For \(X_i=\operatorname{Exp}(x_i/\sqrt\gamma)\) on the complete
chord charts, let \(\Phi_\gamma(X)\) be the positive normalized
product Gaussian vector with its Haar square-root Jacobian divided
out, and with its chart restriction normalized exactly. Its
squared \(x\)-density is the corresponding truncated Gaussian.

Set \(h_{bu}=\operatorname{Exp}((\mathsf F x_b)_u/\sqrt\gamma)\)
at nonroot ports and \(h_{bo}=I\). The moving bridge coordinates are
\[
                    (\Theta_XY)_e=h_s^{-1}Y_eh_t.
\]
At fixed \(X\) this is Haar preserving. Define
\begin{align}
 \widehat S_\gamma f(X,Y)
       &=\Phi_\gamma(X)f(\Theta_XY),\notag\\
 F_{\gamma,t}(Z)
       &=\int|\Phi_\gamma(X)|^2
              M_\times(X,\Theta_X^{-1}Z)^t\,dH^{5B}(X).
 \label{f16v1:gram}
\end{align}
The identities
\[
 \widehat S_\gamma^*M_\times^t\widehat S_\gamma=M_{F_{\gamma,t}},
 \qquad
 \mathcal I_\gamma=M_\times\widehat S_\gamma
                                    M_{F_{\gamma,2}^{-1/2}},
 \qquad \mathcal I_\gamma^*\mathcal I_\gamma=I
\]
are exact on the coarse physical carrier. At each finite regulator
\(F_{\gamma,2}>0\) is bounded below. It is \(F_{\gamma,2}\), not
\(F_{\gamma,1}^2\), that normalizes the magnetic image.

Let \(C_{\rm int}\) and \(D\) be the signed mixed-plaquette incidence
matrices on internal and bridge edges. In the moving frame the
internal linear embedding is \(\bigoplus_bU\), so put
\[
 A=C_{\rm int}\bigoplus_bU,\qquad
 K=AC_BA^{\mathsf T},\quad
 L=C_B^{1/2}A^{\mathsf T}AC_B^{1/2},\quad \rho=1/\sqrt2 .
\]
F15 V99 proves
\[
 0\le K,\quad\|K\|\le\rho<1,\quad
 \Tr K=B\left(\frac3{2\sqrt2}+\frac2{\sqrt{15}}\right)>B.
\]
With \(Z_e=\operatorname{Exp}(y_e/\sqrt\gamma)\),
\begin{equation}
 F_{\gamma,t}(Z)\longrightarrow
 F_t(y)=\det(I+(t/2)K)^{-3/2}
     e^{-(t/4)y^{\mathsf T}
       [D^{\mathsf T}(I+(t/2)K)^{-1}D\otimes I_3]y},
 \label{f16v1:limit}
\end{equation}
uniformly on bounded \(y\)-windows at fixed \(B\).
The proved absolute error is \(O_t(B(1+R^3)/\sqrt\gamma)
+O_t(Be^{-c\gamma})\). It is not a relative bound uniform in volume.

The normalized conditional limit of \(\mathcal I_\gamma\) is the
positive square root of a Gaussian with covariance and mean
\begin{equation}
 C_{\rm p}=(C_B^{-1}+A^{\mathsf T}A)^{-1},\qquad
 m_y=-(C_{\rm p}A^{\mathsf T}D\otimes I_3)y .
 \label{f16v1:posterior}
\end{equation}
The inverse has a block-distance convergent Neumann expansion:
for block sets \(E,F\) at distance \(d\),
\[
 \|P_EC_{\rm p}P_F\|\le\|C\|\rho^d/(1-\rho).
\]
The leading log determinant also has a uniformly summable
per-block closed-walk expansion. These are Gaussian response
statements, not a native non-Gaussian polymer theorem.

F15 V100 counts all internal oscillator levels in this dressed
vector. It proves exponentially small capture by every
\(o(B)\) total-degree cutoff at zero bridge field and an
exponentially accurate cutoff of order
\(B+\|Dy\|^2+\log\delta^{-1}\). Its compact projection handoff
takes \(\gamma\to\infty\) at each fixed \(B\); it does not control
the norm of the full excited-state transfer complement.

\section{New V1 result: exact kinetic compression by conditional affinity}
\label{f16v1:affinity}

For the derivation it is convenient to work first before root Gauss
restriction, on \(L^2(G^{12B})\) and
\(L^2(G^{5B}\times G^{12B})\) in \((X,Z)\) coordinates.
All maps below intertwine the root action, so their exact
operator identities restrict to \eqref{f16v1:physical}.
Define normalized conditional vectors
\begin{equation}
 a_{\gamma,Z}(X)
   =\frac{\Phi_\gamma(X)M_\times(X,\Theta_X^{-1}Z)}
                    {\sqrt{F_{\gamma,2}(Z)}},
 \qquad \int|a_{\gamma,Z}(X)|^2\,dH^{5B}(X)=1,
 \label{f16v1:conditional}
\end{equation}
and their affinity
\[
 \mathcal A_\gamma(Z,Z')
       =\int a_{\gamma,Z}(X)a_{\gamma,Z'}(X)\,dH^{5B}(X).
\]

\begin{theorem}[Exact native bridge kernel on the dressed carrier]
\label{f16v1:exact-kernel}
The compression of the bridge kinetic factor alone is
\begin{equation}
 \bigl(\mathcal I_\gamma^*\mathcal C_\gamma
                      \mathcal I_\gamma\bigr)(Z,Z')
   =\left[\prod_{e\ {\rm bridge}}
             k_\gamma(Z_eZ_e'^{-1})\right]\mathcal A_\gamma(Z,Z').
 \label{eq:f16v1:exact-kernel}
\end{equation}
The affinity is symmetric, positive semidefinite as a Gram kernel,
continuous, and satisfies \(0<\mathcal A_\gamma\le1\) and
\(\mathcal A_\gamma(Z,Z)=1\).
The compressed operator is a positive contraction.
The root gauge invariance is simultaneous in \(Z,Z'\), not
a claim that the affinity is separately invariant in each argument.
\end{theorem}

\begin{proof}
The bridge operator leaves \(X\) fixed. Under
\(Y_e=h_s Z_eh_t^{-1}\) and \(Y_e'=h_s Z_e'h_t^{-1}\),
\[
                  Y_eY_e'^{-1}
                       =h_s Z_eZ_e'^{-1}h_s^{-1}.
\]
Centrality of \(k_\gamma\) and Haar invariance remove both frame
factors. In these coordinates \(\mathcal I_\gamma f=a_{\gamma,Z}f(Z)\).
Integrating the internal coordinate in its compression gives
\eqref{eq:f16v1:exact-kernel}. The Gram assertion and the upper
bound are respectively the squared-norm identity and Cauchy--Schwarz.
Positivity follows also directly from compression of the positive
Wilson convolution. Its norm is at most one.

At finite \(\gamma,B\), the magnetic multiplier is strictly positive
and bounded, \(F_{\gamma,2}\) has a positive lower bound, and
\(\Phi_\gamma\) has norm one. Dominated convergence proves
continuity and strict positivity of the affinity. Equivariance
of the original carrier and kernel proves the root-action
statement and the restriction to invariant functions.
\end{proof}

This identity keeps the whole compact bridge configuration space.
It does not replace \eqref{f16v1:full} by its bridge factor.
In particular the cube kinetic dynamics and their internal
excited-state returns have not been integrated by this theorem.

\subsection{The leading affinity and its exact null directions}

In the next formulas the colour tensor \(I_3\) is suppressed:
when \(R\) acts on a \(36B\)-component bridge vector, it means
\(R\otimes I_3\).
Set
\begin{align}
 R&=D^{\mathsf T}A C_{\rm p}A^{\mathsf T}D
    =D^{\mathsf T}K(I+K)^{-1}D
    =D^{\mathsf T}D-H_{\rm eff},\notag\\
 H_{\rm eff}&=D^{\mathsf T}(I+K)^{-1}D .
 \label{f16v1:R}
\end{align}

\begin{theorem}[Conditional overlap and native local-kernel limit]
\label{f16v1:gaussian-affinity}
For the normalized Gaussian amplitudes \(\psi_y\) with squared
density \eqref{f16v1:posterior},
\begin{equation}
 \langle\psi_y,\psi_z\rangle
             =e^{-(y-z)^{\mathsf T}R(y-z)/8}.
 \label{eq:f16v1:gaussian-affinity}
\end{equation}
At fixed \(B\), \(\mathcal A_\gamma\) in rescaled bridge charts
converges to this expression uniformly on bounded pairs of
bridge windows. After including both square-root Haar chart
densities, the kernel \eqref{eq:f16v1:exact-kernel} converges there to
\begin{equation}
 (2\pi)^{-q/2}
   \exp\left[-\frac{|y-z|^2}{2}
                    -\frac{(y-z)^{\mathsf T}R(y-z)}8\right],
       \qquad q=36B.
 \label{f16v1:native-leading-kernel}
\end{equation}
Thus compression to any fixed bounded coordinate window
converges in Hilbert--Schmidt norm. No global norm or
volume-uniform spectral convergence is asserted.
\end{theorem}

\begin{proof}
Completing the common square of two Gaussian amplitudes with
covariance \(C_{\rm p}\) gives the overlap
\[
       \exp[-(m_y-m_z)^{\mathsf T}C_{\rm p}^{-1}(m_y-m_z)/8].
\]
Substitute \eqref{f16v1:posterior}; Woodbury inversion gives
\eqref{f16v1:R}. F15 V99's normalized-density convergence,
or f15 V100's \(L^2\) amplitude argument, gives the uniform
convergence of the affinities at fixed \(B\).

For completeness, in one bridge chart with Jacobian \(J_\gamma\),
ordinary Taylor expansion and Laplace normalization give
\[
 [J_\gamma(y)J_\gamma(z)]^{1/2}
 k_\gamma(\operatorname{Exp}(y/\sqrt\gamma)
                      \operatorname{Exp}(-z/\sqrt\gamma))
       \longrightarrow(2\pi)^{-3/2}e^{-|y-z|^2/2}
\]
uniformly on compact pairs. Here
\(s(\operatorname{Exp}v)=1-\cos|v|\);
\(\gamma s\) tends to \(|y-z|^2/2\).
The normalization follows by the same rescaling in
\(z_\gamma\); the inequality \(1-\cos r\ge2r^2/\pi^2\)
for \(0\le r\le\pi\) dominates the full chart integral.
The Haar Jacobian cancels to the displayed Euclidean prefactor.
Multiply over the fixed \(12B\) bridges and use
\eqref{eq:f16v1:gaussian-affinity}. Uniform convergence on
a bounded window implies square-integral kernel convergence.
This is a kinematic chart assertion before root averaging;
it does not identify a physical spectral subspace with that window.
\end{proof}

One has
\begin{equation}
 0\le R\le\frac{\rho}{1+\rho}D^{\mathsf T}D.
 \label{f16v1:Rbound}
\end{equation}
The covariance locality from \eqref{f16v1:posterior} and the
finite-range incidence factors give exponentially decaying
block entries of \(R\), uniformly in volume.

There is a stronger exact statement than simply
\(\ker D\subseteq\ker R\). Let \(\mathcal E\) consist of bridge
vectors constant on each set of four parallel bridges.
Then
\begin{equation}
 A^{\mathsf T}D y=0,\qquad R y=0,\qquad
 H_{\rm eff}y=D^{\mathsf T}Dy,\qquad \psi_y=\psi_0
                    \quad (y\in\mathcal E).
 \label{f16v1:common-bridge}
\end{equation}
Indeed a decorated digon takes the difference of two such
parallel values, so its curvature vanishes. A quadrilateral
has no internal edge, so the corresponding row of \(A\) is
zero. This proves the first identity; the others follow.
More generally \(\psi_{y+v}=\psi_y\) for \(v\in\mathcal E\).
The scalar magnetic normalizer still depends on the coarse
quadrilateral curvature; it has not disappeared.

In particular the transverse common-bridge cosine used in
f15 V99 has exactly
\[
       \frac{\langle y,H_{\rm eff}y\rangle}{\|y\|^2}
                           =2\sin^2(\pi/m),
\]
not merely the upper bound established there.
The leading fast conditional vector cannot generate a restoring
term in these coarse common-bridge directions.
This concerns a static linearized Hessian, not a family of
gapless native physical Hamiltonian states.

\section{Finite kinetic time is not fixed by the static normalization}
\label{f16v1:time}

The Gaussian limit allows a complete check of this issue.
Let \(G\) be any positive definite \(q\)-by-\(q\) constant matrix,
and let \(H_t^G\) act on the bridge variable, leaving the internal
variable fixed, with kernel
\[
 h_t^G(y-z)
 =(4\pi t)^{-q/2}\det(G)^{-1/2}
        e^{-(y-z)^{\mathsf T}G^{-1}(y-z)/(4t)}.
\]
The isometry \(\mathcal J f(u,y)=\psi_y(u)f(y)\) compresses it
to \(\mathcal B_t=\mathcal J^*H_t^G\mathcal J\).
The actual one-step local limit \eqref{f16v1:native-leading-kernel}
is the case \(G=I,t=1/2\).

\begin{theorem}[Exact finite-time diffusion change and its composition defect]
\label{f16v1:heat}
For every \(t>0\),
\begin{equation}
 \mathcal B_t=c_t H_t^{G_t},\qquad
 G_t=(G^{-1}+tR/2)^{-1},\qquad
 c_t=\det(I+(t/2)G^{1/2}RG^{1/2})^{-1/2}.
 \label{f16v1:heat-factor}
\end{equation}
Its operator norm on \(L^2(\mathbb R^q)\) is \(c_t\).
If \(R\ne0\), the norm-normalized family
\(c_t^{-1}\mathcal B_t\) is not a semigroup.
\end{theorem}

\begin{proof}
Multiply \(h_t^G(y-z)\) by the affinity
\eqref{eq:f16v1:gaussian-affinity}.
The exponent becomes
\(-(y-z)^{\mathsf T}(G^{-1}+tR/2)(y-z)/(4t)\).
Comparing Gaussian prefactors gives \(c_t\).
The Fourier multiplier is
\(c_t e^{-t p^{\mathsf T}G_t p}\), whose essential supremum is
\(c_t\).

For \(R\ne0\), positivity and inversion show
\(G_t-G_{2t}\) is nonzero positive semidefinite.
Thus for some \(p\),
\[
 e^{-2t p^{\mathsf T}G_t p}
                   \ne e^{-2t p^{\mathsf T}G_{2t}p}.
\]
These are the multipliers of
\((c_t^{-1}\mathcal B_t)^2\) and
\(c_{2t}^{-1}\mathcal B_{2t}\). Continuity in \(p\) makes this
a genuine operator inequality, proving the failure of the
semigroup law.
\end{proof}

The nonzero case actually occurs in the assembly.
For example take a bridge perturbation supported only on
the positive first-direction bridge from fine vertex \((1,0,0)\).
In its incident block, \(A^{\mathsf T}D\) applied to this vector
is nonzero. To see this without a numerical rank test, its
internal-edge source is, up to overall sign, the sum of the
second- and third-direction edges starting at port \((1,0,0)\).
The first/second-direction cube face at third coordinate zero
contains just the former edge, with nonzero signed coefficient.
Thus the source has nonzero circulation. Since
\(\ker U^{\mathsf T}\) is the cut space, its image under
\(U^{\mathsf T}\) cannot vanish. Positivity of \(C_{\rm p}\)
then implies \(R\ne0\).
The exact incidence checker supplies this literal relative-port
example on sides eight and twelve.

The failure is also transparent before taking limits:
\(\mathcal J^*H_s^G H_t^G\mathcal J\) differs from
\(\mathcal J^*H_s^G\mathcal J\mathcal J^*H_t^G\mathcal J\)
by the omitted intermediate internal complement.
A static isometry alone cannot remove that return.

\subsection{The infinitesimal check and its limited physical meaning}

There is a useful independent derivative check.
For bridge indices \(i,j\), real positivity and normalization give
\[
 \langle\psi_y,\partial_i\psi_y\rangle=0,\qquad
 \langle\partial_i\psi_y,\partial_j\psi_y\rangle=R_{ij}/4.
\]
Differentiate the Gaussian mean in
\eqref{f16v1:posterior} to verify the second identity.
For Schwartz \(f\), Fubini's theorem and this zero cross term yield
\begin{equation}
 \int \bigl|\nabla_y(\psi_y f)\bigr|_G^2\,du\,dy
  =\int|\nabla f|_G^2\,dy
                     +\frac14\Tr(GR)\int|f|^2\,dy .
 \label{f16v1:metric}
\end{equation}
Thus the compressed differential quadratic form has a constant
scalar addition, not a bridge-dependent mass potential.
It agrees with the derivative of \eqref{f16v1:heat-factor}:
\[
       \left.\frac d{dt}\log c_t\right|_{t=0}
                           =-\frac14\Tr(GR),\qquad G_0=G.
\]
Exponentiating that compressed differential form is nevertheless
not equal to the finite-time compression in
\eqref{f16v1:heat-factor}. The latter has \(G_t\), not \(G\).
For \(G=I\), the scalar addition is a vacuum normalization of
order at most \(B\), using \eqref{f16v1:Rbound} and the bounded
plaquette incidence. None of these chart statements fixes the
physical lattice time unit or the full interacting transfer gap.

\section{What must be controlled next}
\label{f16v1:next}

The direction is now more precise. The exact native affinity
in \eqref{eq:f16v1:exact-kernel} retains every bridge configuration
and its correct magnetic normalization. Its leading response
has a computable diffusion correction on relative-port directions,
while the fast conditional state is blind to all common-bridge
directions at this order. Repeated ground projection, even after
normalization, discards a definite intermediate return.

The next required calculation is the joint dressed compression
with the cube transfer factors in \eqref{f16v1:full}, retaining
their excited internal bands and the noncommuting moving frame
at both time ends. It must bound the discarded return relative
to the same full native top eigenvalue, uniformly in spatial
volume, or retain it through a controlled memory/Schur term.
A new scalar normalization or another fixed-volume Gaussian
determinant does not supply that estimate.

Beyond it remain nonlinear compact collective-mode coercivity,
bridge-field tail control, an independent physical-time calibration,
and nontrivial interacting continuum reconstruction. These are
requirements of the unchanged objective, not assumptions being
declared solved by the restart.

\section{Verification and preservation}
\label{f16v1:verification}

The diagnostic
\path{scripts/verify_ym_f16_v1_dressed_affinity.py}
checks independent noncommuting matrix completions,
the affinity and derivative metric, finite-time heat precisions,
the composition defect, a direct finite conditional-kernel
compression, and the complete mixed-plaquette incidence on
two actual spatial tori. It checks the common-bridge null
response and an explicit nonzero relative-port response.
It also reruns the inherited f15 V100 diagnostic chain.
All 10,918 new exact fixtures passed, together with the inherited
f15 V100 chain. These finite fixtures supplement proofs, not replace them.
The checker has 5,600 bytes and SHA--256
\begin{center}\scriptsize\ttfamily
CCBCBC50A687E88B74465051A0C7E9B283891911ADA928D910FA5C896B3EFC7B.
\end{center}

F15 V100 is preserved as an archive, with V99's body recoverable
byte for byte. Its source has 2,310,095 bytes and SHA--256
\begin{center}\scriptsize\ttfamily
D9349B9F87D70400E7A140A3BDC38D1BC60F8E3A12F830D8C37693CD6D3C30C5.
\end{center}
This V1 is the sole active main YM source.
Only \texttt{.tex} files are stored in \texttt{latex\_notes};
build products and visual checks are outside it.

\begin{firewall}
The new exact compact result is the conditional-affinity formula
for the dressed bridge kinetic factor. The subsequent Gaussian
overlap, finite-time diffusion change, null directions and
composition defect have proved fixed-volume local limits.
They are not a volume-uniform approximation of the full native
transfer or a positive physical continuum mass-gap theorem.
The full Yang--Mills goal remains unachieved.
\end{firewall}

\section{V2: joint cube dynamics in the dressed transfer}
\label{f16v2:context}

V1 computed the dressed bridge factor while holding the internal
variables fixed. We now include the cube motion at both time ends.
There are two separate corrections: nonroot gauge matching changes
the leading bridge covariance, and compression of the \emph{full}
magnetic sandwich uses the magnetic multiplier twice at each end.
Neither correction is contained in the static conditional affinity.

We first give an exact full-compact kernel. We then integrate its
entire leading internal Gaussian propagator, not just its ground
projection. The resulting two-end Schur matrices have
volume-independent locality and quadratic-form comparisons.
The native limit remains fixed-volume and local in the bridge
coordinates. This does not yet control the discarded native
transfer subspace or the interacting infrared modes.

\subsection{An exact two-ended compact kernel}

Write \(f_v(X)\) for V1's moving port frame, with \(f_o=I\) at
every block root, and put
\[
 \bar X_e=f_s(X)^{-1}X_ef_t(X),\qquad
 \bar X'_e=f_s(X')^{-1}X'_ef_t(X').
\]
This includes tree edges, whose original representatives are \(I\).
As in f15 V99, this is a pointwise gauge rewriting, not a change
of internal integration variables. Denote the exact mixed
multiplier in moving coordinates by
\[
 m_\gamma(X,Z)=M_\times(X,\Theta_X^{-1}Z),
\]
and the internal magnetic multiplier by
\(b_\gamma(X)=\exp[-\gamma S_{\rm int}(X)/2]\).
There are \(7B\) nonroot gauge variables.

\begin{proposition}[Full compact cut kernel with both moving frames]
\label{f16v2:exact-cut}
Before the remaining root Gauss restriction, the cut transfer has
kernel
\begin{align}
 \mathcal K_\gamma(X,Z;X',Z')
 &=b_\gamma(X)b_\gamma(X')\,\mathcal L_\gamma(X,Z;X',Z'),
 \notag\\
 \mathcal L_\gamma
 &=\int_{G^{7B}}
   \prod_{e\ {\rm internal}}
       k_\gamma(r_s^{-1}\bar X_e r_t\bar X_e'^{-1})
   \prod_{e\ {\rm bridge}}
       k_\gamma(r_s^{-1}Z_e r_t Z_e'^{-1})\,dH^{7B}(r).
 \label{f16v2:cut}
\end{align}
Here \(r_o=I\) at block roots. The exact kernel of the full
transfer is \(m_\gamma(X,Z)\mathcal K_\gamma
m_\gamma(X',Z')\). All formulas restrict to the root-invariant
physical space.
\end{proposition}

\begin{proof}
In the rooted tree variables at the two times write \(g'=gh\).
Nonroot Gauss reduction integrates \(h\) and leaves the kinetic
arguments \(h_s^{-1}X_eh_tX_e'^{-1}\) on internal edges and
\(h_s^{-1}Y_eh_tY_e'^{-1}\) on bridges.
This is the same exact Fourier-free Haar reduction used in
f15 V97; no boundary representation has been discarded.

Set \(h_v=f_v(X)r_vf_v(X')^{-1}\).
For fixed endpoints this is a product of Haar translations.
The internal argument becomes, up to conjugation by \(f_s(X')\),
\(r_s^{-1}\bar X_e r_t\bar X_e'^{-1}\).
On a bridge, substitute \(Y_e=f_s(X)Z_ef_t(X)^{-1}\) and its
primed counterpart. The same conjugation gives
\(r_s^{-1}Z_er_tZ_e'^{-1}\). Centrality of \(k_\gamma\) removes
the conjugations. The internal and mixed magnetic traces are
unchanged by the pointwise gauge rewriting.
This proves the kernel and the full sandwich.
\end{proof}

Let \(\mathcal I_\gamma\) be V1's exact normalized dressed
isometry. Distinguish
\[
 \mathcal B_{1,\gamma}=\mathcal I_\gamma^*
                  T_{\rm cut,\gamma}\mathcal I_\gamma,\qquad
 \mathcal B_{2,\gamma}=\mathcal I_\gamma^*
                  T_{\rm full,\gamma}\mathcal I_\gamma.
\]
Their exact kernels, for \(\ell=1,2\), are
\begin{equation}
 \begin{aligned}
 \mathcal B_{\ell,\gamma}(Z,Z')
 &=\frac{1}{\sqrt{F_{\gamma,2}(Z)F_{\gamma,2}(Z')}}\\
 &\quad\times\int \Phi_\gamma(X)\Phi_\gamma(X')\,
      m_\gamma(X,Z)^\ell m_\gamma(X',Z')^\ell
      \mathcal K_\gamma(X,Z;X',Z')\,dH^{5B}(X)dH^{5B}(X').
 \end{aligned}
 \label{f16v2:exact-compression}
\end{equation}
In particular, \(\ell=2\) is necessary for the full transfer:
\[
 \mathcal I_\gamma^*T_{\rm full,\gamma}\mathcal I_\gamma
 =M_{F_{\gamma,2}^{-1/2}}\widehat S_\gamma^*
      M_\times^2T_{\rm cut,\gamma}M_\times^2\widehat S_\gamma
                                      M_{F_{\gamma,2}^{-1/2}}.
\]
Using \(\ell=1\) instead would compress a different operator.

\subsection{The kinetic covariance produced by the moving nonroot vertices}

Use block-diagonal copies \(\mathsf B_B,\Sigma_B,\mathsf M_B,
\mathsf H_B,U_B\) of the single-cube matrices in V1 and f15 V97,
where \(\Sigma=(\mathsf B^{\mathsf T}\mathsf B)^{-1}\).
Let \(E\) be the \(12B\)-by-\(7B\) bridge endpoint incidence:
its row is \(z_t-z_s\), with each root coordinate set to zero.
Define
\begin{equation}
 G=I+E\Sigma_BE^{\mathsf T},\qquad
 L_h=\mathsf B_B^{\mathsf T}\mathsf B_B+E^{\mathsf T}E.
 \label{f16v2:G}
\end{equation}
The letter \(G\) in these matrix formulas denotes a covariance;
the group remains \(\SU(2)\).

\begin{theorem}[Exact leading kinetic square and a uniform covariance bound]
\label{f16v2:kinetic-square}
For one colour component, \(\delta x=x-x'\), \(\delta y=y-y'\),
and \(z\in\mathbb R^{7B}\),
\begin{align}
 &\|U_B\delta x+\mathsf B_Bz\|^2+\|\delta y+Ez\|^2
 \notag\\
 &\quad=\delta x^{\mathsf T}\mathsf M_B^{-1}\delta x
       +\delta y^{\mathsf T}G^{-1}\delta y
       +(z+L_h^{-1}E^{\mathsf T}\delta y)^{\mathsf T}
                    L_h(z+L_h^{-1}E^{\mathsf T}\delta y).
 \label{f16v2:completed-kinetic}
\end{align}
Moreover
\begin{equation}
  \det L_h=(\det\mathsf M)^B\det G,\qquad
                         I\le G\le30I.
 \label{f16v2:Gbounds}
\end{equation}
The matrix \(G\) couples only bridges incident to a common block.
\end{theorem}

\begin{proof}
The balanced internal embedding obeys
\(U_B^{\mathsf T}\mathsf B_B=0\) and
\(U_B^{\mathsf T}U_B=\mathsf M_B^{-1}\).
There is therefore no internal--port cross term.
Complete the \(z\)-square. Woodbury inversion gives
\[
 I-E L_h^{-1}E^{\mathsf T}
                    =(I+E\Sigma_BE^{\mathsf T})^{-1}=G^{-1}.
\]
The determinant lemma and the cube equality
\(\det(\mathsf B^{\mathsf T}\mathsf B)=\det\mathsf M\)
give the determinant statement.

Each row of \(E\) has absolute sum at most two. Each nonroot
port is incident to three bridges, so each column has absolute
sum three. Thus \(\|E\|^2\le6\).
The diagonal of the rooted inverse \(\Sigma\) consists of the
effective resistances from the root. For the cube these are
\(7/12\) at the three adjacent vertices, \(3/4\) at the three
face-opposite vertices and \(5/6\) at the opposite vertex.
They are the exact resistance values verified and derived
in f15 V98. Consequently
\[
 \|\Sigma_B\|=\|\Sigma\|\le\Tr\Sigma
                    =3(7/12)+3(3/4)+5/6=29/6.
\]
This proves \(G\le I+6(29/6)I=30I\).
Block diagonality of \(\Sigma_B\) gives the support assertion.
\end{proof}

The bound \(30\) is a convenient pre-root-quotient bound,
not a replacement for f15 V98's sharper physical electric
form comparison. The new point is that the bridge covariance
is \(G\), not \(I\). Fixing all temporal port variables to zero
would miss it.

For a positive \(d\)-by-\(d\) path covariance \(Q\), abbreviate
the three-colour Gaussian density by
\[
 \varphi_Q(v)
  =(2\pi)^{-3d/2}\det(Q)^{-3/2}
             e^{-v^{\mathsf T}(Q^{-1}\otimes I_3)v/2}.
\]
Write \(b_0(x)=e^{-x^{\mathsf T}(\mathsf H_B\otimes I_3)x/4}\).
The complete leading cut kernel is
\begin{equation}
 \mathcal K_0(x,y;x',y')
       =b_0(x)b_0(x')\,
           \varphi_{\mathsf M_B}(x-x')\varphi_G(y-y').
 \label{f16v2:leading-cut}
\end{equation}
The full leading kernel has in addition
\(e^{-|Ax+Dy|^2/4}\) at each end.

\begin{proposition}[Fixed-volume native convergence with all internal ends integrated]
\label{f16v2:native-limit}
At each fixed \(B\), after including all square-root Haar
chart densities, the kernel \eqref{f16v2:cut} converges uniformly
on bounded \(x,x',y,y'\) windows to \eqref{f16v2:leading-cut}.
The kernels \eqref{f16v2:exact-compression} converge uniformly
on bounded \(y,y'\) windows to the result of integrating the
complete Gaussian kernel \eqref{f16v2:leading-cut} with its
displayed endpoint weights and normalizers.
There is no internal excitation cutoff in this assertion.
It also gives Hilbert--Schmidt convergence after restricting
both bridge ends to a fixed bounded coordinate window.
\end{proposition}

\begin{proof}
Put \(r_v=\operatorname{Exp}(z_v/\sqrt\gamma)\).
The first-order kinetic words in \eqref{f16v2:cut} are exactly
\(U_B(x-x')+\mathsf B_Bz\) and \(y-y'+Ez\).
Their actions have leading exponent one half of the sum of
the two squared norms. Gaussian integration and
\eqref{f16v2:completed-kinetic} give the covariance in
\eqref{f16v2:leading-cut}; \eqref{f16v2:Gbounds} fixes its
normalizing determinant. The internal plaquette multiplier
converges to \(b_0\).

Here is the compact-integral justification.
When the endpoint charts tend to identity, the internal tree
kinetic words force every nonroot \(r_v\) to identity: each
tree has a fixed root. On the compact group the sum of their
squared geodesic distances controls the distances of all
\(r_v\), up to an endpoint error bounded by \(C/\gamma\)
on a fixed endpoint window. This follows successively along
the seven tree edges, by the bi-invariant triangle inequality
and Cauchy--Schwarz along each path. Since
\(1-\cos\theta\ge2\theta^2/\pi^2\) for \(0\le\theta\le\pi\),
the rescaled integrand is dominated by
\(C\exp(-c\sum_v|z_v|^2)\), with constants at fixed \(B\)
and window. The complement of any fixed small group
neighbourhood is exponentially small. Taylor expansion there,
followed by dominated convergence, proves the local kernel limit.

For the integration over all internal \(X,X'\), a separate
uniform bound avoids a hidden chart restriction.
Keep only the \(7B\) tree factors of the kinetic product and
bound the other \(17B\) factors by \(\|k_\gamma\|_\infty\).
Successive Haar integration of leaves gives one for the tree
product. The full integral is therefore at most
\(\|k_\gamma\|_\infty^{17B}\).
The normalization in \eqref{f16v1:kinetic} gives
\(\|k_\gamma\|_\infty\le C\gamma^{3/2}\) for \(\gamma\ge1\).
Each chart Jacobian is at most \(C\gamma^{-3/2}\).
The two half-densities on the \(17B\) reduced coordinates
cancel this power. Thus the cut chart kernel is bounded by
a constant depending on \(B\), uniformly throughout the
internal charts; magnetic multipliers are at most one.

The transformed \(\Phi_\gamma\) is exactly a normalized
truncation of the product Gaussian ground amplitude.
It therefore gives an integrable dominating product
\(C_B\Omega(x)\Omega(x')\).
On a fixed bridge window \(F_{\gamma,2}\) is bounded below
uniformly for sufficiently large \(\gamma\), by V1's
fixed-volume limit. Internal Gaussian tails are now uniform
in \(\gamma\). Apply the local limit on bounded internal
windows and then send those windows to infinity.
This proves uniform convergence of the compressed kernels.
The final Hilbert--Schmidt assertion follows on any bounded
bridge window. All constants may deteriorate with \(B\).
\end{proof}

\subsection{The two-end Schur formula retaining the full internal propagator}

The product cube Gaussian ground amplitude \(\Omega(x)\)
has squared covariance \(C_B\) from \eqref{f16v1:covariance}.
Write \(N=\mathsf M_B^{-1}\) and
\begin{equation}
 Q_0=\tfrac12(C_B^{-1}+\mathsf H_B),\qquad
 V_{\ell,+}=Q_0+\tfrac\ell2 A^{\mathsf T}A,\qquad
 V_{\ell,-}=V_{\ell,+}+2N,\qquad \ell=0,1,2.
 \label{f16v2:V}
\end{equation}
All are strictly positive. Define
\begin{align}
 P_{\ell,\pm}&=D^{\mathsf T}A V_{\ell,\pm}^{-1}A^{\mathsf T}D,
 \notag\\
 B_{\ell,\pm}
       &=\ell D^{\mathsf T}D-H_{\rm eff}
                              -\tfrac{\ell^2}{2}P_{\ell,\pm}.
 \label{f16v2:B}
\end{align}
Colour tensors are suppressed: each path matrix acts with \(I_3\)
on the actual bridge vectors.

Let \(r(q)=(1+q/2+\sqrt{q+q^2/4})^{-1}\), and put
\[
                   \lambda_\Box^0=r(4)^{9/2}r(6)^3.
\]
This is the top value of the entire single-cube Gaussian
transfer, not its physical mass. It is the
\href{https://dlmf.nist.gov/18.18.E28}{Mehler value}
already derived in f15 V97, with nine coordinates of parameter
four and six coordinates of parameter six.
Set
\begin{equation}
 \kappa_\ell
 =\frac{(\lambda_\Box^0)^B}{F_2(0)}
   \left[
    \frac{\det V_{0,+}\det V_{0,-}}
         {\det V_{\ell,+}\det V_{\ell,-}}
   \right]^{3/2},\qquad \ell=1,2.
 \label{f16v2:kappa}
\end{equation}

\begin{theorem}[Complete Gaussian dressed cut and full compressions]
\label{f16v2:schur}
The fixed-volume limits in Proposition~\ref{f16v2:native-limit} are
\begin{align}
 \mathcal B_{\ell,0}(y,y')
 &=\kappa_\ell\,\varphi_G(y-y') \notag\\
 &\quad\cdot
 \exp\left[-\tfrac18(y+y')^{\mathsf T}B_{\ell,+}(y+y')
           -\tfrac18(y-y')^{\mathsf T}B_{\ell,-}(y-y')\right],
       \qquad \ell=1,2.
 \label{f16v2:schur-kernel}
\end{align}
In particular \(\ell=2\) includes the entire internal Gaussian
propagator of the full magnetic transfer between the dressed ends.
No intermediate internal ground projection appears.
\end{theorem}

\begin{proof}
At one end the integrand of \eqref{f16v2:exact-compression}
tends to
\[
 \frac{\Omega(x)
           e^{-\ell|Ax+Dy|^2/4}b_0(x)}{\sqrt{F_2(y)}} .
\]
The two-end internal precision and source are
\[
 W_\ell=
 \begin{pmatrix}V_{\ell,+}+N&-N\\-N&V_{\ell,+}+N\end{pmatrix},
 \qquad
 \eta_\ell=-\frac\ell2
          \begin{pmatrix}A^{\mathsf T}Dy\\A^{\mathsf T}Dy'\end{pmatrix}.
\]
Changing to \((x+x')/\sqrt2,(x-x')/\sqrt2\) gives the two
precisions \(V_{\ell,+},V_{\ell,-}\). Gaussian integration
therefore contributes determinant
\((\det V_{\ell,+}\det V_{\ell,-})^{-3/2}\)
and exponent
\[
 \frac{\ell^2}{16}\left[
   (y+y')^{\mathsf T}P_{\ell,+}(y+y')
     +(y-y')^{\mathsf T}P_{\ell,-}(y-y')\right].
\]
The remaining endpoint exponent is
\[
 -\frac\ell4(y^{\mathsf T}D^{\mathsf T}Dy
                   +y'^{\mathsf T}D^{\mathsf T}Dy')
 +\frac14(y^{\mathsf T}H_{\rm eff}y
                   +y'^{\mathsf T}H_{\rm eff}y').
\]
The second term comes from \(F_2(y)^{-1/2}F_2(y')^{-1/2}\).
Combining these expressions gives exactly \eqref{f16v2:B}.

To fix the scalar without discarding any determinant, set the
endpoint magnetic power to zero in the internal integral.
It is \(\langle\Omega,K_{\Box,0}^{\otimes B}\Omega\rangle
=(\lambda_\Box^0)^B\). Dividing the general determinant by
this one and retaining \(F_2(0)^{-1}\) gives
\eqref{f16v2:kappa}. This proves the formula.
\end{proof}

This is a calculation of one-step compression, not a claim that
its spectrum equals the spectrum of \(T_{\rm full,\gamma}\).
V1's intermediate-complement warning remains in force.

\subsection{Uniform locality of the retained two-time fast response}

The inverses in \eqref{f16v2:V} are uniformly local even for the
full magnetic power \(\ell=2\).
For either sign put
\[
 Q_+=Q_0,\quad Q_-=Q_0+2N,\qquad
 Z_\pm=Q_\pm^{-1/2}A^{\mathsf T}A Q_\pm^{-1/2}.
\]
Since \(Q_\pm\ge C_B^{-1}/2\),
\[
                 0\le Z_\pm,\qquad \|Z_\pm\|\le2\rho .
\]
These matrices have only nearest-block off-diagonal entries.
Let \(c=1+2\rho\), \(\theta=1-c^{-1}<1\), and
\(T_\pm=I-(I+Z_\pm)/c\). Then \(0\le T_\pm\le\theta I\) and
\begin{equation}
 V_{2,\pm}^{-1}
   =c^{-1}Q_\pm^{-1/2}
                     \sum_{j=0}^{\infty}T_\pm^jQ_\pm^{-1/2}.
 \label{f16v2:local-inverse}
\end{equation}
For block sets \(E_0,F_0\) separated by distance \(d\), this gives
\[
 \|P_{E_0}V_{2,\pm}^{-1}P_{F_0}\|
                    \le\|Q_\pm^{-1}\|\theta^d.
\]
Indeed powers of degree less than \(d\) cannot connect the sets,
and \(c^{-1}/(1-\theta)=1\).
No volume or set-cardinality factor is needed.
Consequently \(P_{2,\pm}\) and \(B_{2,\pm}\) have exponentially
decaying block entries, with the fixed range shift of \(A,D\).
This controls the complete leading two-time fast response,
not the non-Gaussian native remainder.

\subsection{The full effective quadratic transfer and its surviving soft modes}

Write \(B_\pm=B_{2,\pm}\).
The lower bound here is relative to the actual coarse curl,
not relative to the identity.

\begin{theorem}[Volume-independent curvature and kinetic comparisons]
\label{f16v2:comparison}
One has
\begin{equation}
 \frac{D^{\mathsf T}D}{(1+\rho)(1+2\rho)}
       \le B_+\le B_-
       \le\frac{1+2\rho}{1+\rho}D^{\mathsf T}D .
 \label{f16v2:Bbounds}
\end{equation}
Define
\[
       C_{\rm eff}
       =\left[G^{-1}+\tfrac14(B_--B_+)\right]^{-1}.
\]
Then
\begin{equation}
 \frac{I}{1+16\rho/(1+2\rho)}
       \le C_{\rm eff}\le30I ,
 \label{f16v2:Ceffbounds}
\end{equation}
and, with
\(\widetilde\kappa=\kappa_2(\det C_{\rm eff}/\det G)^{3/2}\),
\begin{equation}
 \mathcal B_{2,0}(y,y')
   =\widetilde\kappa\,
      e^{-y^{\mathsf T}B_+y/4}
       \varphi_{C_{\rm eff}}(y-y')
      e^{-y'^{\mathsf T}B_+y'/4}.
 \label{f16v2:effective}
\end{equation}
\end{theorem}

\begin{proof}
Since \(V_{2,+}\ge C_B^{-1}/2+A^{\mathsf T}A\),
\[
 A V_{2,+}^{-1}A^{\mathsf T}
              \le2K(I+2K)^{-1}.
\]
The commuting functions of \(K\) obey the exact identity
\[
 2I-(I+K)^{-1}-4K(I+2K)^{-1}
                    =(I+K)^{-1}(I+2K)^{-1}.
\]
Substitute in \(B_+=2D^{\mathsf T}D-H_{\rm eff}-2P_{2,+}\)
and use \(0\le K\le\rho I\) for the lower bound.
The upper bound follows from \(P_{2,-}\ge0\) and
\(H_{\rm eff}\ge(1+\rho)^{-1}D^{\mathsf T}D\).
Inversion of \(V_{2,-}\ge V_{2,+}\) gives \(B_-\ge B_+\).

Furthermore
\[
 0\le B_--B_+
   =2D^{\mathsf T}A(V_{2,+}^{-1}-V_{2,-}^{-1})A^{\mathsf T}D
   \le\frac{4\rho}{1+2\rho}D^{\mathsf T}D.
\]
Every bridge belongs to four plaquettes; each mixed row has at
most four bridge entries. Thus \(\|D\|^2\le16\).
Together with \(I\le G\le30I\), this proves
\eqref{f16v2:Ceffbounds}. Finally split the exponent in
\eqref{f16v2:schur-kernel} into the two endpoint \(B_+\)
terms and the difference term \(B_--B_+\).
The Gaussian prefactor changes by the displayed determinant
ratio, proving \eqref{f16v2:effective}.
\end{proof}

These comparisons prove \(\ker B_+=\ker D\).
For the common-bridge subspace \(\mathcal E\) of V1,
\(A^{\mathsf T}Dy=0\) and \(H_{\rm eff}y=D^{\mathsf T}Dy\).
Thus
\[
 B_+y=B_-y=D^{\mathsf T}Dy,\qquad
            (B_--B_+)y=0\quad(y\in\mathcal E).
\]
The full leading cube propagation has not generated a constant
restoring term for those coarse collective modes.

For clarity, even the positive oscillator frequencies of
\eqref{f16v2:effective} have a soft finite-volume edge.
Let \(\nu_{\min}^+\) be the least positive eigenvalue of
\(C_{\rm eff}^{1/2}B_+C_{\rm eff}^{1/2}\).
Then
\begin{equation}
                     \nu_{\min}^+\le60\sin^2(\pi/m).
 \label{f16v2:soft}
\end{equation}
To prove it, use V1's common-bridge transverse cosine \(v\),
orthogonal in the Euclidean metric to \(\ker D\).
Its numerator is \(\langle v,D^{\mathsf T}Dv\rangle\).
Choose the representative \(v+w\), \(w\in\ker D\), orthogonal
to that kernel in the \(C_{\rm eff}^{-1}\) metric.
Its numerator is unchanged, and its denominator is at least
\(\|v+w\|^2/30\ge\|v\|^2/30\).
The generalized minimum principle and V1's exact quotient
\(2\sin^2(\pi/m)\) prove \eqref{f16v2:soft}.

The zero directions include residual linearized gauge directions
and torons. The full Euclidean Gaussian operator has free
diffusion in these directions, not a normalizable isolated vacuum.
If those directions are removed to discuss just its positive
oscillator factors, a parameter \(\nu>0\) has Mehler ratio
\((1+\nu/2+\sqrt{\nu+\nu^2/4})^{-1}\).
Equation~\eqref{f16v2:soft} concerns this specified Gaussian
factorization. It is neither a native physical gaplessness
theorem nor a choice of physical time scaling.

\subsection{The remaining full-operator return and verification}

The exact compact kernel and the full Gaussian Schur response
now include both time ends and all internal propagation within
one step. The missing spectral implication is still substantive:
a one-step compression does not bound the full complement.
With \(P_\gamma=\mathcal I_\gamma\mathcal I_\gamma^*\),
the actual two-step return is
\[
 \mathcal I_\gamma^*T_{\rm full,\gamma}
       (I-P_\gamma)T_{\rm full,\gamma}\mathcal I_\gamma .
\]
It is nonnegative, but no sufficiently small relative bound
uniform in spatial volume has been proved. The native bridge
tails and the nonlinear common-bridge dynamics must also be
kept under the same normalization. The fixed-volume domination
in Proposition~\ref{f16v2:native-limit} is not such a bound.

The next upstream target is to retain this return as a
two-time or Schur/memory operator and test a volume-uniform
relative comparison, rather than identify the compressed
Gaussian spectrum with the full transfer spectrum.
Independent physical-time calibration and interacting
four-dimensional continuum reconstruction remain open.

The diagnostic
\begin{center}\small
\path{scripts/verify_ym_f16_v2_joint_dressed_transfer.py}
\end{center}
checks noncommuting two-ended compact words, the actual
two-cube kinetic completion and determinant, full-torus port
incidences, the independent two-end Gaussian precision and
plus/minus Schur formulas, the matrix identity behind
\eqref{f16v2:Bbounds}, and the endpoint magnetic-square
distinction. Finite checks supplement the proofs; they are not
a continuum spectral certificate. All 4,582 new exact fixtures passed,
together with the inherited V1 chain. The diagnostic uses exact
elimination for the larger determinants, rather than the inherited
small-matrix cofactor routine. It has 8,154 bytes and SHA--256
\begin{center}\scriptsize\ttfamily
95087B3BF6B4AA9517C734E8A35B690B9074EF3D9355AC7125B2A22340C59DCF.
\end{center}

The predecessor V1 has 25,601 bytes and SHA--256
\begin{center}\scriptsize\ttfamily
A04BB24BC54769C31EF36F591282F6ABCF828DA25920C369DCB9DA9F1EF27097.
\end{center}
Its body is preserved byte for byte; removing this append and
restoring the title recovers that source. All 167 previously
protected files remain unchanged. The last companion
checkpoint is f15 V100 and the next is this lineage's V5.
Only \texttt{.tex} files belong in \texttt{latex\_notes};
build and visual-check products are kept outside that folder.

\begin{firewall}
V2 gives an exact full-compact joint kernel, a fixed-volume
native dressed-compression limit with all internal ends
integrated, and explicit volume-independent comparisons for
its complete leading two-time Gaussian response.
No small volume-uniform full-operator return, positive
physical continuum mass gap, or interacting continuum
Yang--Mills construction is proved. The goal remains unachieved.
\end{firewall}

\section{V3: extensive dressed-line leakage and a retained fast history}
\label{f16v3:context}

V2 computed one full step between the normalized dressed ends.
We now ask what happens before inserting another dressed projection.
The answer is not a small correction to the static conditional line:
even the complete leading Gaussian transfer changes its fast
covariance by an extensive amount. We prove a quantitative
pre-root-restriction obstruction to a small relative return.
We then integrate the fast variables along an arbitrary bridge
history without intermediate projection. Its covariance and
quadratic memory have volume- and time-length-independent decay.

This is a change in the next analytic target, not a mass-gap claim.
F15 V99--V100 concerned the magnetic image and its old oscillator
occupations; V1 concerned bridge-only affinity; V2 concerned the
joint one-step compression. The new calculation propagates the
dressed image by the full Gaussian transfer and retains arbitrarily
many fast time slices. The archive search also found spatial
Gaussian blocking memory in f11 V32--V33 and native observable
covariance estimates in f15 V61; neither is the conditional
fast-history calculation below.

\subsection{The propagated conditional vector is not the static dressed vector}

All statements in this subsection concern V2's leading Gaussian
model before the remaining root Gauss restriction.
Write \(\mathsf T_0\) for its full transfer on the internal and
bridge Euclidean coordinates, and
\[
 (\mathcal J_0 f)(x,y)=\psi_y(x)f(y),\qquad
 \psi_y(x)=\frac{\Omega(x)e^{-|Ax+Dy|^2/4}}{\sqrt{F_2(y)}} .
\]
The vector \(\psi_y\) has norm one in the internal coordinates.
Each displayed path matrix is tensored with \(I_3\) when acting
on the three colour components. Put
\begin{align}
 N&=\mathsf M_B^{-1},&
 S&=\mathsf H_B+A^{\mathsf T}A,&
 \Lambda&=A^{\mathsf T}D,\notag\\
 E_0&=N+\tfrac12(C_B^{-1}+\mathsf H_B),&
 E&=E_0+A^{\mathsf T}A,&
 P&=\tfrac12(C_B^{-1}+A^{\mathsf T}A).
 \label{f16v3:definitions}
\end{align}
Here \(E\) is V2's two-end diagonal block \(V_{2,+}+N\).
The precision in the amplitude \(\psi_z\), as opposed to its
squared probability density, is \(P\):
\[
 \psi_z(x)=c(z)\exp[-x^{\mathsf T}Px/2-x^{\mathsf T}\Lambda z/2],
                      \qquad c(z)>0.
\]
The sources and the positive scalar depend on the bridge field.

For fixed initial bridge coordinate \(y\) and final coordinate \(z\),
define the fast output feature
\[
 q_{z,y}(x)=\int \mathsf T_0(x,z;x',y)\psi_y(x')\,dx'.
\]
This definition has no intermediate fast projection.

\begin{proposition}[Full-step covariance mismatch]
\label{f16v3:propagation}
Set
\begin{equation}
 \Delta_f=N(E_0^{-1}-E^{-1})N\ge0.
 \label{f16v3:delta}
\end{equation}
There is a positive scalar \(c(z,y)\) such that
\begin{equation}
 q_{z,y}(x)=c(z,y)
 \exp\left[
   -\tfrac12x^{\mathsf T}(P+\Delta_f)x
   -x^{\mathsf T}\bigl(\tfrac12\Lambda z+NE^{-1}\Lambda y\bigr)
 \right].
 \label{f16v3:feature}
\end{equation}
Consequently the squared normalized output feature has covariance
\([2(P+\Delta_f)]^{-1}\otimes I_3\), independent of \(y,z\).
It is not the covariance \((2P)^{-1}\otimes I_3\) of the static
dressed line unless \(A=0\).
\end{proposition}

\begin{proof}
The full kernel is the product of its internal and bridge heat
densities and the two endpoint factors
\[
 \exp[-(x^{\mathsf T}\mathsf H_Bx+|Ax+Dz|^2)/4].
\]
After multiplication by \(\psi_y(x')\), the input precision is \(E\)
and its linear source is \(Nx-\Lambda y\).
Gaussian integration of \(x'\) leaves output precision
\[
 Q=N+\tfrac12S-NE^{-1}N
\]
and linear source \(-\Lambda z/2-NE^{-1}\Lambda y\).
The cube ground-state equation gives the exact identity
\begin{equation}
 N+\tfrac12(\mathsf H_B-C_B^{-1})=NE_0^{-1}N.
 \label{f16v3:riccati}
\end{equation}
For completeness, conjugate by \(\mathsf M_B^{1/2}\).
On a cube mode with parameter \(t=4\) or \(6\), the two sides
become
\[
 1+t/2-\sqrt{t+t^2/4}
       =\bigl(1+t/2+\sqrt{t+t^2/4}\bigr)^{-1};
\]
their product identity follows by squaring the radical.
This proves \eqref{f16v3:riccati}, hence \(Q=P+\Delta_f\).
Inversion reverses \(E\ge E_0\), so \(\Delta_f\ge0\).
If \(\Delta_f=0\), invertibility of \(N,E,E_0\) implies
\(A^{\mathsf T}A=0\). The covariance assertion follows by squaring
the amplitude and normalizing its positive Gaussian.
\end{proof}

\subsection{A volume-extensive obstruction to a small static return}

The covariance mismatch gives a uniform conditional overlap
bound, not just a nonzero example.

\begin{theorem}[Dressed-line retention after a full step]
\label{f16v3:retention}
For the actual parity-cube geometry with \(B=m^3\), define
\[
 \eta_B(z,y)=
       \frac{|\langle\psi_z,q_{z,y}\rangle|^2}
            {\|q_{z,y}\|_2^2}.
\]
Then, for every \(y,z\),
\begin{equation}
 0<\eta_B(z,y)\le
 \eta_B^{\rm cov}:=
 \frac{\det(I+Z)^{3/2}}{\det(I+Z/2)^3}
 \le e^{-c_*B}\le e^{-B/100000},
 \qquad
 c_*=\frac{50}{147}\frac{(25/1728)^2}{5}>10^{-5}.
 \label{f16v3:retention-bound}
\end{equation}
Here \(Z\) is the relative covariance-precision change defined
in the proof. The determinants retain all \(5B\) path modes,
with the displayed power accounting for all three colours.
\end{theorem}

\begin{proof}
Use bars to denote congruence by \(N^{-1/2}=\mathsf M_B^{1/2}\);
in particular
\[
 \overline E_0=N^{-1/2}E_0N^{-1/2},\quad
 W=N^{-1/2}A^{\mathsf T}AN^{-1/2},\quad
 \overline P=N^{-1/2}PN^{-1/2}.
\]
Then
\[
 \overline\Delta_f
   =\overline E_0^{-1}-(\overline E_0+W)^{-1},\qquad
 Z=\overline P^{-1/2}\overline\Delta_f\,\overline P^{-1/2}.
\]
The eigenvalues of \(Z\) are the generalized eigenvalues of
\((\Delta_f,P)\); no commuting-matrix assumption is being made.

The inherited cube spectrum and covariance give
\[
 5I\le\overline E_0\le8I,\qquad 2I\le\overline P\le6I.
\]
Indeed the exact bounds for \(\overline E_0\) are
\(3+2\sqrt2\) and \(4+\sqrt{15}\).
For the second assertion, the spectrum of the whitened
\(C_B^{-1}\) lies between \(4\sqrt2\) and \(2\sqrt{15}\), and
\(0\le W\le4I\).
This last inequality follows from
\(A=C_{\rm int}U_B\), \(\|C_{\rm int}\|\le2\), and
\(U_B^{\mathsf T}U_B=N\).
Moreover the exact diagonal Gram blocks
\((A^{\mathsf T}A)_{bb}=2N_{bb}\) from f15 V99 give
\[
                           \Tr W=10B.
\]
These are bounds for the actual geometry, not assumptions
on independently chosen toy matrices.

Woodbury inversion, or continuous extension if \(W\) is singular,
gives
\[
 \overline\Delta_f
 =\overline E_0^{-1}W^{1/2}
   (I+W^{1/2}\overline E_0^{-1}W^{1/2})^{-1}
                       W^{1/2}\overline E_0^{-1}
 \ge\frac59\,\overline E_0^{-1}W\overline E_0^{-1}.
\]
Taking traces and using \(\overline E_0\le8I\) yields
\[
 \Tr\overline\Delta_f\ge\frac{25}{288}B,\qquad
 \Tr Z\ge\frac{25}{1728}B,\qquad 0\le Z\le\frac1{10}I.
\]
For the last estimate use
\(\overline\Delta_f\le\overline E_0^{-1}\le I/5\) and
\(\overline P\ge2I\).
Cauchy--Schwarz on the \(5B\) eigenvalues gives
\[
                   \Tr Z^2\ge\frac{(25/1728)^2}{5}\,B.
\]

The squared overlap of two normalized Gaussian amplitudes with
precisions \(P\) and \(P+\Delta_f\) is at most the value when
their means coincide. Completing the square shows that unequal
means multiply this value by \(e^{-a}\) for a nonnegative
quadratic form \(a\). The coincident-mean determinant is precisely
\(\eta_B^{\rm cov}\) in \eqref{f16v3:retention-bound}; congruence changes
do not change it. For an eigenvalue \(0\le z\le1/10\),
\[
 \frac32\log\frac{(1+z/2)^2}{1+z}
 =\frac32\log\left(1+\frac{z^2}{4(1+z)}\right)
 \ge\frac{3z^2}{8(1+z/2)^2}
 \ge\frac{50}{147}z^2 .
\]
We used \(\log(1+u)\ge u/(1+u)\), which follows by integrating
its nonnegative derivative difference from zero.
Summation proves the stated exponential bound.
\end{proof}

Let
\[
 \mathsf B_0=\mathcal J_0^*\mathsf T_0\mathcal J_0,\quad
 \mathsf A_2=\mathcal J_0^*\mathsf T_0^2\mathcal J_0,\quad
 \mathsf R_0=\mathsf A_2-\mathsf B_0^2
 =\mathcal J_0^*\mathsf T_0(I-\mathcal J_0\mathcal J_0^*)
                              \mathsf T_0\mathcal J_0\ge0.
\]
Their diagonal kernels satisfy
\begin{align}
 \mathsf A_2(y,y)&=\int\|q_{z,y}\|_2^2\,dz,\notag\\
 \mathsf B_0^2(y,y)&=\int|\langle\psi_z,q_{z,y}\rangle|^2\,dz
                 \le e^{-c_*B}\mathsf A_2(y,y).
 \label{f16v3:diagonal}
\end{align}
All these diagonal integrals are finite. In fact the internal
cut transfer has norm \((\lambda_\Box^0)^B\); its magnetic
multipliers have norm at most one. Hence
\[
 \|q_{z,y}\|_2
       \le(\lambda_\Box^0)^B\varphi_G(z-y),
\]
whose square is integrable in \(z\), uniformly in \(y\).

For any bounded bridge window \(U\) of positive measure,
the resulting operators localized at their input are
Hilbert--Schmidt and
\begin{equation}
 \|(I-\mathcal J_0\mathcal J_0^*)\mathsf T_0
                      \mathcal J_0\,1_U\|_{\rm HS}^2
 \ge(1-e^{-c_*B})
              \|\mathsf T_0\mathcal J_0\,1_U\|_{\rm HS}^2.
 \label{f16v3:HS-loss}
\end{equation}
This follows by integrating \eqref{f16v3:diagonal}, or by
orthogonal decomposition of each output feature.
In particular a quadratic-form inequality
\(\mathsf R_0\le\varepsilon\mathsf A_2\) on the entire
pre-root-restriction coarse space requires
\(\varepsilon\ge1-e^{-c_*B}\). To see this, sandwich by \(1_U\)
and take the finite trace; the denominator is strictly positive.
Thus no fixed \(\varepsilon<1\) works for all volumes in this model.

This conclusion is stronger than a nonzero projection defect,
but its scope is important. It does not assert pointwise
positivity of every off-diagonal return kernel, nor does it
turn a diagonal comparison into a reverse operator inequality.
It does not apply automatically after physical Gauss
restriction, to a low-energy subspace, or to the native compact
transfer. In particular, V2's fixed-window one-step convergence
alone does not transfer this full two-step statement to that
native operator.

\subsection{Exact elimination along an arbitrary bridge history}

We now avoid intermediate projection altogether. Let \(n\ge1\),
fix the bridge history \(\mathbf y=(y_0,\ldots,y_n)\), and
integrate the internal coordinates \(x_0,\ldots,x_n\) in the
kernel of \(\mathcal J_0^*\mathsf T_0^n\mathcal J_0\).
Set \(E_{\rm i}=S+2N\), and define the \((n+1)\)-by-\((n+1)\)
block precision
\begin{equation}
 (W_n)_{00}=(W_n)_{nn}=E,\qquad
 (W_n)_{ii}=E_{\rm i}\ (0<i<n),\qquad
 (W_n)_{i,i+1}=(W_n)_{i+1,i}=-N,
 \label{f16v3:history-precision}
\end{equation}
with all other time blocks zero.
Its time blocks have \(5B\) path coordinates.
Let \(\mathcal L\mathbf y=(\Lambda y_0,\ldots,\Lambda y_n)\) and
\[
 a_n=\frac1{F_2(0)}
       \left[\det C_B\,(\det\mathsf M_B)^n\det W_n\right]^{-3/2}.
\]

\begin{theorem}[Complete fast-history kernel]
\label{f16v3:history}
After integration of every internal time slice, the exact
Gaussian history weight is
\begin{align}
 \mathcal W_n(\mathbf y)
 &=a_n\prod_{i=0}^{n-1}\varphi_G(y_i-y_{i+1})
       \exp\bigl(\mathcal E_n(\mathbf y)\bigr),\notag\\
 \mathcal E_n(\mathbf y)
 &=-\frac12\sum_{i=0}^n y_i^{\mathsf T}D^{\mathsf T}Dy_i
   +\frac14\bigl(y_0^{\mathsf T}H_{\rm eff}y_0
                         +y_n^{\mathsf T}H_{\rm eff}y_n\bigr)
   +\frac12(\mathcal L\mathbf y)^{\mathsf T}
                           W_n^{-1}(\mathcal L\mathbf y).
 \label{f16v3:history-weight}
\end{align}
The kernel of \(\mathcal J_0^*\mathsf T_0^n\mathcal J_0\) is
\(\int\mathcal W_n(\mathbf y)\,dy_1\cdots dy_{n-1}\).
No internal state or intervening bridge configuration has
been discarded. The determinant in \(a_n\) is part of this
identity, not a freely chosen spectral normalization.
\end{theorem}

\begin{proof}
Each interior time slice receives two magnetic endpoint
multipliers, whereas each external slice receives one
transfer multiplier and one multiplier from its dressed
end vector. Thus the mixed magnetic power is two at
\emph{every} slice. Each linear source is consequently
\(-\Lambda y_i\).
The internal magnetic power is two at interior slices;
at the two ends there is one such factor and one
\(\Omega\) amplitude. The kinetic differences contribute
\(\sum_{i=0}^{n-1}(x_i-x_{i+1})^{\mathsf T}N(x_i-x_{i+1})\).
These observations give exactly \eqref{f16v3:history-precision}.
Its quadratic form can also be written as
\begin{align*}
 \mathbf x^{\mathsf T}W_n\mathbf x
 ={}&x_0^{\mathsf T}(E-N)x_0+x_n^{\mathsf T}(E-N)x_n\\
 &+\sum_{i=1}^{n-1}x_i^{\mathsf T}Sx_i
   +\sum_{i=0}^{n-1}(x_i-x_{i+1})^{\mathsf T}N(x_i-x_{i+1}),
\end{align*}
which proves strict positivity.
Gaussian integration contributes the last term in
\(\mathcal E_n\) and \(\det(W_n)^{-3/2}\).
The scalar \((2\pi)^{15B(n+1)/2}\) from integration cancels
the \(n\) internal heat prefactors and the two ground-amplitude
prefactors. Their remaining factors are
\((\det\mathsf M_B)^{-3n/2}(\det C_B)^{-3/2}\).
Finally the two factors \(F_2^{-1/2}\) contribute \(F_2(0)^{-1}\)
and the displayed external \(H_{\rm eff}\) term.
This proves the history formula. Positivity permits the
iterated integrations; their interpretation as the stated
bounded-operator kernel follows from composition of the
positive Gaussian transfer kernels.
\end{proof}

For \(n=1\), the plus/minus transform has precisions
\(E-N=V_{2,+}\) and \(E+N=V_{2,-}\), so this is exactly V2,
including its scalar. For every \(n\) one may compute the
determinant without dropping its boundary terms by
\[
 R_0=E,\qquad R_i=(W_n)_{ii}-NR_{i-1}^{-1}N,\qquad
                   \det W_n=\prod_{i=0}^n\det R_i .
\]
All pivots are strictly positive by block Gaussian elimination.
The order of the matrix products matters.

\subsection{Uniform temporal and space--time locality of the fast memory}

Write \(\mathbb N=I_{n+1}\otimes N\) and
\(\widehat W_n=\mathbb N^{-1/2}W_n\mathbb N^{-1/2}\).
The graph for space--time locality has a vertex for each
(time slice, spatial block), with the nearest-time and
nearest-spatial-block edges.

\begin{theorem}[Length-independent conditional fast response]
\label{f16v3:locality}
For every \(B,n\),
\begin{equation}
                        4I\le\widehat W_n\le15I.
 \label{f16v3:uniform-spectrum}
\end{equation}
For sets \(\mathcal E,\mathcal F\) at graph distance \(d\),
\begin{equation}
 \|P_{\mathcal E}\widehat W_n^{-1}P_{\mathcal F}\|
                         \le\frac14(11/15)^d .
 \label{f16v3:spacetime}
\end{equation}
A sharper purely temporal bound is
\begin{equation}
 \|N^{1/2}(W_n^{-1})_{ij}N^{1/2}\|
       \le\frac{q^{|i-j|}}{p_*-2},
 \qquad p_*=3+2\sqrt2,\quad q=2/p_*<1.
 \label{f16v3:temporal}
\end{equation}
These are conditional fast-covariance bounds for arbitrary
bridge histories, not decay of the complete physical
Yang--Mills correlations.
\end{theorem}

\begin{proof}
Whitening the quadratic-form decomposition in the preceding
proof gives a time-path Laplacian, with norm at most four.
The interior on-site term is the whitened
\(\mathsf H_B+A^{\mathsf T}A\), between \(4I\) and \(10I\).
The external on-site term is the whitened \(E-N\), between
\((2+2\sqrt2)I\) and \((7+\sqrt{15})I<11I\).
This proves \eqref{f16v3:uniform-spectrum}.
All off-site entries of \(\widehat W_n\) join adjacent
space--time blocks: \(A^{\mathsf T}A\) has only nearest-block
entries, and the whitening is spatially block diagonal.
Thus \(T_n=I-\widehat W_n/15\) satisfies
\(0\le T_n\le(11/15)I\), with the same finite range.
The convergent expansion
\[
             \widehat W_n^{-1}=\frac1{15}\sum_{k\ge0}T_n^k
\]
has no connection between those sets before degree \(d\).
The geometric tail is \((11/15)^d/4\), proving
\eqref{f16v3:spacetime} without a set-cardinality factor.

For the temporal estimate, keep entire spatial slices as blocks.
The diagonal blocks of \(\widehat W_n\) are at least \(p_*I\):
the external ones are \(\overline E_0+W\), and the interior
ones are at least \(6I\). Its temporal off-diagonal blocks are
\(-I\). If \(\mathcal D_n\) is its time-block diagonal,
write \(\widehat W_n=\mathcal D_n-\mathcal A_n\), where
\(\mathcal A_n\) is the path adjacency with identity blocks.
Then
\(\|\mathcal D_n^{-1/2}\mathcal A_n\mathcal D_n^{-1/2}\|
\le2/p_*=q\).
This matrix need not be positive, but its norm-convergent
Neumann series is valid. A term of degree less than
\(|i-j|\) has zero \(ij\) block.
Multiplication by the two diagonal inverse square roots and
summation give
\(p_*^{-1}q^{|i-j|}/(1-q)=q^{|i-j|}/(p_*-2)\).
Undoing the whitening proves \eqref{f16v3:temporal}.
\end{proof}

There is also a curvature-relative truncation estimate.
Let \(\mathcal Q_n(\mathbf y)\) denote the last term in
\(\mathcal E_n\), and let \(\mathcal Q_n^{[r]}\) keep only
its ordered pairs with \(|i-j|\le r\), for an integer \(r\ge0\).
Since \(\|N^{-1/2}A^{\mathsf T}\|\le2\), the temporal bound yields
\begin{equation}
 |\mathcal Q_n(\mathbf y)-\mathcal Q_n^{[r]}(\mathbf y)|
 \le \frac{4}{p_*-2}\frac{q^{r+1}}{1-q}
                               \sum_{i=0}^n\|Dy_i\|^2 .
 \label{f16v3:memory-tail}
\end{equation}
Indeed each ordered matrix element is bounded by
\(4q^{|i-j|}\|Dy_i\|\|Dy_j\|/(p_*-2)\).
The omitted scalar matrix has row sum at most
\(2q^{r+1}/(1-q)\); its quadratic form is bounded by that
row sum times \(\sum_i\|Dy_i\|^2\). Include the factor
\(1/2\) in \(\mathcal Q_n\) to obtain the displayed constant.
No spatial volume or time-length factor multiplies this
relative curvature budget. Truncation is an estimate on the
quadratic action, not automatically an operator-norm or
positivity-preserving approximation.

For histories taking values in V1's common-bridge subspace,
\(\Lambda y_i=0\) at every time. The entire fast source
response \(\mathcal Q_n\), not merely its long-time tail,
then vanishes. The determinant \(a_n\) is independent of
the bridge history. This statement concerns the fixed-history
weight; integrating other bridge modes is a further operation.
It does not create a constant restoring term or settle the
native dynamics of those common modes.

\subsection{Direction of the next native estimate and verification}

The exact native counterpart of the return is still
\[
 \mathcal I_\gamma^*T_{\rm full,\gamma}
        (I-\mathcal I_\gamma\mathcal I_\gamma^*)
                         T_{\rm full,\gamma}\mathcal I_\gamma.
\]
V3 shows why asking for a small relative return on the entire
static dressed carrier is not a promising unrestricted
Gaussian target: \eqref{f16v3:HS-loss} goes the opposite way.
A physical low-energy restriction would require its own
estimate; it is not supplied by this obstruction.

The constructive replacement is the unprojected fast-history
integral. Its Gaussian precision remains uniformly coercive
even though the retained bridge dynamics has soft directions.
The next useful native theorem must control the non-Gaussian
conditional fast integral, its determinant or pressure
normalization, and its bridge-field tails in a norm compatible
with the physical Gauss restriction and the same full transfer
top value. The compact action is not globally convex in a
single chart, and V2's constants depend on volume. Consequently
\eqref{f16v3:memory-tail} is not yet a native uniform cluster or
spectral estimate. Physical-time calibration and an interacting
continuum construction remain additional open requirements.

The new diagnostic is
\begin{center}\small
\path{scripts/verify_ym_f16_v3_fast_history.py}.
\end{center}
Its 2,075 exact rational fixtures check the noncommuting
propagated precision and source, unequal-mean overlap
completion, full-history action and determinant recursion,
one-step agreement with V2, finite-range propagation, and
length-independent memory tails. They also check the rational
constants in the extensive leakage bound. The inherited V2
chain passed. Finite fixtures supplement the proofs and do
not certify a physical mass gap.
The checker has 7,716 bytes and SHA--256
\begin{center}\scriptsize\ttfamily
811A66019E7A8B4A78D6CF8C1F8F6B2882D318BBA5D2B57D6E4D5C274275EA02.
\end{center}

The predecessor V2 has 48,007 bytes and SHA--256
\begin{center}\scriptsize\ttfamily
D8BC72F2B7E92E94C2B8E3AF848D9580563E5F9302CA13E8843909C42AB80F6F.
\end{center}
Its body is preserved byte for byte: remove this append and
restore the title to recover it. The 168 previously protected
files are unchanged. The last companion checkpoint is f15 V100;
the next is f16 V5. Only \texttt{.tex} files belong in
\texttt{latex\_notes}; build products remain elsewhere.

\begin{firewall}
V3 proves extensive static dressed-line loss and a conditional,
unprojected fast-history formula with uniform Gaussian
space--time locality and curvature-relative temporal truncation.
The leakage obstruction is before physical Gauss restriction.
No native volume-uniform memory bound, nontrivial interacting
continuum construction, or positive physical Yang--Mills
mass gap is established. The full goal remains unachieved.
\end{firewall}

\end{document}
